\newpage{}
\section{From Edges to Nodes}
\label{sec:edges-to-nodes}

\subsection{Balance Equation}
\label{subsec:balance-equation}

With the graph convention from Section~\ref{sec:edges-and-nodes}, a node view
becomes available when an edge price is interpreted as a ratio of node
valuations.  Let $\eta_u\in\Rpos$ denote the node valuation of node
$u\in\mathcal{V}$.  For edge $e=(u,v)\in\mathcal{E}$, the idealized edge price
is
\[
  q_e
  =
  \frac{\eta_u}{\eta_v}
  \in\Rpos.
\]
In log coordinates,
\[
  \ell_e := \log q_e = \xi_u-\xi_v,
  \qquad
  \ell_e\in\R,
  \qquad
  \xi_u := \log\eta_u\in\R.
\]
Therefore a signed scalar edge return satisfies
\[
  r_e
  =
  \Delta\ell_e
  =
  y_u-y_v,
  \qquad
  r_e\in\R,
  \qquad
  y_u := \Delta\xi_u\in\R.
\]
This is the basic balance equation.  The edge forecast is a relation; the node
state is a potential whose differences explain that relation.

\subsection{Best-Fit Node Recovery}
\label{subsec:best-fit-node-recovery}

In practice the edge readouts may be noisy, masked, or internally inconsistent.
The node returns are then recovered as the best global fit to the selected edge
returns.  Let $\mathcal{E}_{\mathrm{use}}\subseteq\mathcal{E}$ denote the
usable edges for one forecast step and one scalar observer.  To keep notation
light, the equations below write sums over $\mathcal{E}$; one may read this as
summing over $\mathcal{E}_{\mathrm{use}}$ after masking invalid edges.

For each usable edge, let $\widehat{r}_e\in\R$ be the observed or predicted edge
return and let $\gamma_e\in\Rpos$ be its confidence weight.  Recover the node
vector $\vect{y}\in\R^N$ by
\[
  \min_{\vect{y}\in\R^N}
  \sum_{e=(u,v)\in\mathcal{E}}
  \gamma_e\cdot
  \left(
    \widehat{r}_e-(y_u-y_v)
  \right)^2
  \qquad\text{subject to}\qquad
  y_g=0.
\]
The gauge node $g\in\mathcal{V}$ fixes the additive ambiguity: only differences
of node values matter.  If the usable graph is disconnected, one gauge is
needed for each connected component; otherwise the recovery is not unique on
isolated components.

\subsection{Incidence-Matrix Form}
\label{subsec:incidence-matrix-form}

Let $\mat{A}\in\R^{L\times N}$ be the directed incidence matrix of
$\mathcal{G}$.  For edge $e=(u,v)$, the corresponding incidence entries are
\[
  A_{e\,u}=1,
  \qquad
  A_{e\,v}=-1,
\]
and zeros elsewhere.  Let $\vect{\widehat{r}}\in\R^L$ collect the edge
observations and let
\[
  \mat{\Gamma}
  =
  \diag(\gamma_1,\ldots,\gamma_L)
  \in\R^{L\times L}.
\]
Then node recovery is the weighted least-squares problem
\[
  \min_{\vect{y}\in\R^N}
  \lVert
    \mat{\Gamma}^{1/2}
    (\vect{\widehat{r}}-\mat{A}\vect{y})
  \rVert_2^2,
  \qquad
  y_g=0.
\]
Its normal equations are
\[
  \mat{A}^{\top}\mat{\Gamma}\mat{A}\vect{y}
  =
  \mat{A}^{\top}\mat{\Gamma}\vect{\widehat{r}},
  \qquad
  y_g=0.
\]
The matrix
\[
  \mat{L}_{\gamma}
  :=
  \mat{A}^{\top}\mat{\Gamma}\mat{A}
  \in\R^{N\times N}
\]
is the confidence-weighted graph Laplacian.  Without gauge fixing it has the
usual constant-vector null direction on each connected component; the gauge
constraint removes that ambiguity for the solve.

\subsection{Recovery Residuals}
\label{subsec:recovery-residuals}

After recovery, the remaining edge mismatch is
\[
  \vect{\eps}_r
  =
  \vect{\widehat{r}}
  -
  \mat{A}\widehat{\vect{y}}
  \in\R^L.
\]
This residual measures the part of the edge-return field that cannot be
explained by one global node-potential vector.  A weighted scalar diagnostic is
\[
  \mathcal{Q}_r
  =
  \left\|
    \mat{\Gamma}^{1/2}\vect{\eps}_r
  \right\|_2^2
  \in\Rnonneg.
\]
Large residuals indicate disagreement among edge forecasts, poor confidence
weights, or cycle inconsistency in the active market graph.

\subsection{Covariance-Aware Recovery}
\label{subsec:covariance-aware-recovery}

If a full edge covariance estimate $\mat{C}_r\in\R^{L\times L}$ is available for
the chosen observer, then the natural generalization is
\[
  \min_{\vect{y}\in\R^N}
  \left(
    \vect{\widehat{r}}-\mat{A}\vect{y}
  \right)^{\top}
  \mat{C}_r^{-1}
  \left(
    \vect{\widehat{r}}-\mat{A}\vect{y}
  \right),
  \qquad
  y_g=0.
\]
This is the Mahalanobis-weighted version of the same balance equation.  It does
not introduce a new node model; it only measures edge mismatch using the edge
uncertainty structure instead of independent scalar weights.

When $\mat{C}_r$ is singular or nearly singular, the implementation should use
a regularized inverse or a pseudoinverse on the usable edge subspace.  This is a
numerical protection around the same mathematical object, not a change in the
meaning of the recovery.

\subsection{Node Covariance and Geometry}
\label{subsec:node-covariance-and-geometry}

When the covariance-aware system is well posed, the recovered node covariance is
read after applying the same gauge fixing used in the least-squares solve.  If
$\mat{A}_g$ denotes the incidence matrix in the free node coordinates after the
gauge coordinate has been removed, then
\[
  \mat{C}_y^{(g)}
  \approx
  \left(
    \mat{A}_g^{\top}\mat{C}_r^{-1}\mat{A}_g
  \right)^{-1} 
  \in \R^{D_p}.
\]
Equivalently, one may keep all $N$ coordinates and understand the covariance on
the quotient space where additive constants are identified, using the
corresponding constrained inverse or pseudoinverse.  The important point is that
the null gauge direction is not an economic uncertainty mode.

This covariance gives a natural local geometry on recovered node states.  For
two gauge-fixed node configurations $\vect{y},\vect{y}'$, the Mahalanobis
distance is
\[
  d(\vect{y},\vect{y}')^2
  =
  (\vect{y}-\vect{y}')^{\top}
  \mat{C}_y^{-1}
  (\vect{y}-\vect{y}')
  \in\Rnonneg.
\]
This distance measures separation in units of node uncertainty, not in raw
coordinate scale.

\subsection{Modes}
\label{subsec:modes}

If a covariance matrix is available on the recovered node space, diagonalize the
gauge-fixed covariance in the free coordinates as
\[
  \mat{C}_y^{(g)}
  =
  \mat{V}\Lambda\mat{V}^{\top},
  \qquad
  \mat{V}\in\R^{(N-1)\times(N-1)},
  \qquad
  \Lambda\in\R^{(N-1)\times(N-1)},
\]
when one gauge node has been removed.  In quotient-space notation the same modes
may be lifted back to the $N$ node coordinates, with the gauge direction excluded.
The eigenvectors are the principal modes of collective node motion, and the
eigenvalues measure how much variance each mode carries.  Keeping only a few
dominant modes is an optional compression step, not a mandatory reduction.

\calloutbox{End of the Forecasting Interface.}{\\%
  At this point the document has defined the active graph, normalized data,
  representation, edge-level predictive distributions, observer readouts, and
  the optional lift from scalar edge observations to node coordinates.  Trading
  policy, execution rules, and portfolio decision logic remain outside this
  notation layer.
}
